\documentclass[modern]{aastex701}

\usepackage{amsmath}

\shorttitle{Citizen Photometry}
\shortauthors{Author et al.}

\begin{document}

\title{Citizen Photometry: A Local-First Desktop Pipeline for Turning Amateur Astrophotography into Shareable Time-Domain Measurements}

\correspondingauthor{Author Name Placeholder}
\email{author@example.com}

\author{Author Name Placeholder}
\email{author@example.com}
\affiliation{Affiliation Placeholder}

\begin{abstract}
Amateur astrophotographers collectively acquire a large volume of astronomical imaging data, yet most of those observations never enter scientific workflows. The limiting factor is often not data acquisition itself, but the difficulty of converting heterogeneous local image folders into trustworthy photometric products without specialized domain knowledge, custom scripts, or cloud-based data sharing. We present \emph{Citizen Photometry}, a local-first desktop application that guides users from folders of FITS or XISF images to inspectable time-domain outputs while keeping raw data on the observer's own machine. The software combines automated file scanning, world-coordinate-system validation with optional \texttt{astrometry.net} solving, catalog-assisted source selection using Gaia DR3, VSX, and the NASA Exoplanet Archive, aperture photometry, differential magnitude estimation, interactive review of measurements and annotated images, and export of science-oriented report bundles including AAVSO-oriented products. The project is motivated by a citizen-science use case: enabling non-professional observers to contribute scientifically useful measurements from existing image archives and ongoing observing programs without requiring them to surrender local control of their underlying image data. This manuscript presents the problem setting, design requirements, software architecture, and current workflow in a form suitable for later quantitative evaluation on representative variable-star case studies.
\end{abstract}

\keywords{Astronomy software --- Time domain astronomy --- Photometry --- Variable stars --- Citizen science --- Astronomical databases}

\section{Introduction}\label{sec:introduction}

Amateur astronomy produces a large amount of observational data outside formal survey programs. In particular, astrophotographers routinely collect repeated imaging sequences of the night sky for hours at a time, often over many nights or observing seasons. Those datasets can contain scientifically useful time-domain information, especially for known variable stars and other sources that benefit from long-baseline monitoring. However, only a small fraction of that observational effort is converted into scientific measurements suitable for review, aggregation, or follow-up.

The core bottleneck is procedural rather than purely instrumental. Many observers already own equipment capable of producing useful photometric sequences, but the reduction path from local image folders to a scientifically interpretable light curve remains difficult. Existing workflows often assume familiarity with calibration and catalog concepts, command-line tooling, file-format quirks, and manual quality control. For many non-professional observers, that barrier is high enough that potentially useful data remain archived only as personal image collections.

At the same time, many observers are reluctant to upload their raw image libraries to third-party platforms. In practice this may reflect data-volume constraints, privacy preferences, ownership concerns, or simple unwillingness to expose a full archive when only derived scientific products need to be shared. A useful citizen-science tool in this setting should therefore reduce technical barriers to analysis while preserving local control over source data.

Citizen Photometry was developed to address that gap. The software is a desktop application that processes local astronomical image sequences into inspectable photometric outputs while making automation the default and human oversight readily available. The project is aimed at known-target workflows rather than blind all-sky discovery, with current emphasis on differential photometry and practical report generation.

\section{Motivation And Problem Setting}\label{sec:motivation}

The motivating problem for this work is that amateur astrophotographers generate a potentially valuable observational resource, but much of it remains outside scientific use because the reduction pathway is too specialized, too fragmented, or too dependent on expert intervention.

Several constraints recur in this setting.

\begin{enumerate}
\item Users may have substantial observing experience but limited formal background in scientific photometric reduction.
\item Scientific software stacks often assume scripting competence and comfort with astronomical catalogs, metadata conventions, and calibration choices.
\item Upload-based workflows conflict with a common preference to keep the original image archive private and local.
\item Even when motivated to contribute, users may not know which targets are scientifically useful, which outputs are worth exporting, or how to assess whether an automated result is credible.
\end{enumerate}

The consequence is a mismatch between observational effort and scientific reuse. If that mismatch can be reduced, then existing astrophotography archives could become a meaningful source of citizen-science measurements, especially for variable-star monitoring, eclipse timing, and long-baseline follow-up. The present project is an attempt to build software infrastructure for that transition.

\section{Contributions}\label{sec:contributions}

This paper makes four primary contributions.

\begin{enumerate}
\item It frames amateur astrophotography archives as a local, underutilized source of time-domain observations that can be made more scientifically accessible through software rather than through new instrumentation alone.
\item It presents a local-first desktop architecture that keeps raw files on the user's machine while producing shareable scientific outputs such as measurement tables, light curves, annotated images, and report bundles.
\item It describes an integrated workflow that combines file discovery, WCS handling, catalog lookup, target preview, aperture photometry, differential analysis, interactive inspection, and export in a single application.
\item It establishes a practical software foundation for future evaluation on real observing campaigns, particularly for variable-star follow-up and other known-target use cases.
\end{enumerate}

The present manuscript is intentionally scoped as a software-and-motivation paper. It does not yet claim a broad quantitative scientific validation across many targets. Instead, it defines the problem, the operational model, and the current implementation in a form suitable for later case-study and benchmark work.

\section{Related Work}\label{sec:related}

Citizen Photometry sits at the intersection of several existing traditions in astronomy software: professional and semi-professional photometric reduction pipelines, astronomical catalog services, citizen-science contribution frameworks, and general-purpose astronomy Python infrastructure.

At the infrastructure layer, the current implementation depends on the Python scientific-astronomy ecosystem, including \texttt{Astropy} and \texttt{astroquery} for file handling, coordinate logic, catalog access, and related utilities \citep{astropy2013,astropy2018,ginsburg2019}. The program also relies on public catalog resources including Gaia DR3 \citep{gaia-dr3}, the AAVSO International Variable Star Index (VSX) \citep{vsx}, and the NASA Exoplanet Archive \citep{nasa-exoplanet-archive}. When image headers do not contain usable astrometric solutions, the workflow can optionally use \texttt{astrometry.net} \citep{astrometrynet}.

From a workflow perspective, many existing tools already support high-quality photometric analysis, but they often assume a higher level of scripting fluency, reduction expertise, or willingness to assemble a multi-tool pipeline. Citizen Photometry is not intended to replace those tools in expert settings. Instead, it focuses on a different operating point: providing a guided, inspectable, desktop workflow for non-professional observers who already have image data but are unlikely to build a bespoke analysis stack.

The project is also adjacent to citizen-science systems that aggregate volunteer effort around scientific classification or monitoring tasks. The distinction here is that Citizen Photometry targets the conversion of personally held observational data into sharable scientific products. In that sense, it is not simply a participation portal, but a local data-reduction and triage environment for potential contributors.

\section{Design Requirements}\label{sec:requirements}

The design was guided by four practical requirements.

\begin{description}
\item[Local-first processing.] Raw images should remain on the observer's own machine by default, and sharing should occur mainly through derived outputs.
\item[Accessible automation.] Common operations such as scanning, metadata extraction, WCS validation, and photometric measurement should not require scripting.
\item[Inspectability.] Automated outputs should remain auditable through tables, annotated images, light curves, and diagnostic notes.
\item[Scientific usefulness.] Exports should be oriented toward downstream review and community science workflows rather than only toward visual presentation.
\end{description}

These requirements led to a desktop application rather than a cloud platform, and to a pipeline that exposes intermediate artifacts instead of hiding them behind an opaque batch process.

\section{Software Overview}\label{sec:overview}

Citizen Photometry is currently a Windows-first desktop application written in Python. It accepts FITS and XISF images, supports opening a project root, a dedicated \texttt{Files} directory, or a single object directory, and centers its main workflow on known targets rather than blind transient discovery. The application persists workspace settings locally, caches catalog responses and run history, and provides both automated processing and manual inspection tools.

At a high level, the software performs the following steps.

\begin{enumerate}
\item Scan local image folders and extract available metadata from image headers and filenames.
\item Validate existing WCS information and optionally request plate solutions when needed.
\item Build a representative field description and query catalogs for nearby Gaia sources, known VSX variables, and exoplanet host entries.
\item Offer a preview-driven source-selection stage for choosing which targets to process.
\item Measure fluxes through aperture photometry using fixed or FWHM-scaled apertures.
\item Select comparison stars and compute differential photometric products.
\item Present file-level, source-level, measurement-level, and image-level diagnostics in the desktop interface.
\item Export report bundles, tabular outputs, annotated images, light-curve plots, and AAVSO-oriented files.
\end{enumerate}

This integrated workflow is meant to collapse what would otherwise be a loosely connected chain of tools into a single guided application.

\section{Pipeline And Methods}\label{sec:methods}

\subsection{Input Handling And Metadata Extraction}\label{subsec:inputs}

The software scans local directory trees for supported astronomical image formats, currently including FITS and XISF. The intended organization is a project-level \texttt{Files} directory containing one subdirectory per observed object, but the application also accepts narrower entry points such as a single object folder. During scanning, the program reads available header metadata and can infer additional structure from filenames when headers are incomplete.

This early normalization stage is operationally important because astrophotography archives are often heterogeneous. File naming, metadata completeness, and astrometric state vary substantially across observers and across software used upstream in acquisition or stacking.

\subsection{WCS Validation And Optional Astrometric Solving}\label{subsec:wcs}

Scientific reuse of image data requires reliable sky-coordinate context. Citizen Photometry therefore validates existing WCS metadata before attempting source selection and catalog matching. If a usable solution is not available, the pipeline can request an external plate solve through \texttt{astrometry.net} \citep{astrometrynet}. This is optional and only needed when local headers do not already contain sufficient WCS information.

The current implementation also includes practical tolerance for real-world amateur data products, including partially specified WCS headers and image formats that may encode astrometric information outside a conventional FITS keyword set. This is a software-engineering choice driven by the reality of user-supplied archives rather than idealized survey data.

\subsection{Catalog-Assisted Source Selection}\label{subsec:catalogs}

Once a solved field is available, the program queries multiple public resources to contextualize the image sequence. Gaia DR3 provides the broad field-star context \citep{gaia-dr3}; VSX supplies known variable-star identifications \citep{vsx}; and the NASA Exoplanet Archive contributes cataloged exoplanet-host entries \citep{nasa-exoplanet-archive}. Catalog responses are cached locally so that repeated inspection of the same field does not require repeated remote queries.

The catalog stage is used both for context and for triage. Rather than requiring the user to begin from an empty field, the software can present a preview of likely targets and allow selection before full processing. This keeps the workflow focused on scientifically interesting sources while still leaving room for manual choices.

\subsection{Photometric Measurement Workflow}\label{subsec:photometry}

The current science workflow is centered on differential photometry. The pipeline measures source fluxes through aperture photometry, supports fixed-radius or FWHM-scaled aperture profiles, and selects comparison stars for relative brightness estimation. Resulting measurements can then be organized into per-source time series and rendered as inspectable light curves.

The emphasis here is practical usability rather than maximal methodological novelty. The software is intended to help users obtain a credible, reviewable photometric product from their local image archive with fewer manual steps, not to introduce a fundamentally new photometric estimator. Its contribution lies in packaging the end-to-end workflow in a form that reduces operational friction for non-specialist observers.

\subsection{Interactive Inspection And Manual Oversight}\label{subsec:inspection}

An important design principle is that automation should remain inspectable. The desktop interface therefore exposes multiple synchronized views of the processing result, including file-level summaries, source-level summaries, measurement tables, annotated images, and light-curve panels. These views allow users to assess which files contributed to a result, how a source appears in the image plane, whether outliers are present, and whether comparison-star choices appear reasonable.

The application also supports manual interaction alongside automated analysis. This matters because citizen-science quality control often depends on user judgment: rejecting problematic frames, checking whether a measurement trend is plausible, or inspecting whether a target and comparison stars were handled sensibly in crowded or noisy data.

\subsection{Export Products}\label{subsec:exports}

Citizen Photometry exports multiple output types intended for review and downstream use. Current products include measurement CSV files, light-curve CSV files, summary JSON files, plot images, annotated images, accepted and rejected science-observation tables, provenance-oriented manifests, and AAVSO-oriented text and preflight outputs. This export model reflects the project's local-first philosophy: raw source images can remain private while derived products are prepared for sharing and further scientific assessment.

The export bundle is intentionally conservative. It aims to produce useful, structured downstream artifacts now while leaving room for later refinement of publication-grade provenance, transformation, and submission workflows.

\section{Local-First Contribution Model}\label{sec:localfirst}

The software is explicitly designed around a local-first contribution model. Users keep their image data on local storage and generate sharable outputs from that local archive. This design addresses a practical sociotechnical barrier in citizen science: some observers are willing to contribute results, but not willing to upload entire datasets, expose private observing archives, or surrender control of imagery that may have personal or competitive value.

By separating raw data retention from derived-output sharing, Citizen Photometry attempts to make scientific participation more acceptable to that user group. In this model, the shared unit is not the original image library, but the reduced product: a measurement table, diagnostic plot set, annotated frame set, or report bundle that another party can inspect and review.

\section{Scope And Limitations}\label{sec:limitations}

The current system is deliberately constrained. It focuses on desktop execution, local file processing, and known-target workflows, especially variable stars. It is not yet a blind transient-discovery pipeline, and it should not be described as a full replacement for expert reduction environments. Likewise, the present paper does not yet provide a broad quantitative validation campaign across many observing setups.

Several evaluation tasks remain for future work: measuring usability for non-specialists, benchmarking workflow compression relative to multi-tool alternatives, validating results against representative known targets, and estimating the scale of dormant amateur image archives that could realistically be converted into shareable scientific products.

\section{Future Work}\label{sec:future}

The most important next step is to pair the software description in this paper with case studies on real objects. A strong follow-on paper would include representative variable-star targets, comparisons against trusted reference reductions, error-characterization discussion, and a clearer estimate of how much latent observing time in the astrophotography community could be mobilized through this style of tool.

Additional software directions include improved submission-oriented exports, broader support for manual review workflows, and expansion toward candidate-event triage while retaining human validation. The key question is not only whether the software works technically, but whether it can measurably increase scientifically useful participation from observers who would otherwise remain outside formal reduction pipelines.

\section{Conclusion}\label{sec:conclusion}

Citizen Photometry is motivated by a simple mismatch: a large community already produces astronomical image data, but only a small portion of that work becomes structured scientific output. A local-first desktop pipeline can reduce that mismatch by lowering operational complexity, preserving user control over raw files, and producing reviewable photometric outputs suitable for citizen-science contribution. In that sense, the project is not merely a convenience application for image analysis; it is an attempt to build practical infrastructure between amateur observing practice and time-domain scientific reuse.

\begin{acknowledgments}
This manuscript is currently a project draft. Author list, institutional affiliations, acknowledgments, and quantitative evaluation remain to be finalized.
\end{acknowledgments}

\software{Python, Astropy, astroquery, Matplotlib, NumPy, pandas, photutils, PySide6, pyqtgraph, Citizen Photometry}

\bibliographystyle{aasjournalv7}
\bibliography{references}

\end{document}